%! Unit Launch: Study Design — Unit 1, Lesson 1.0 — Warm-Up
% ONE PAGE, blank and key. Three items at 12pt (the stats-model size), \workrowsep 50pt, item
% spacing 22pt, so each item gets real handwriting room and still fits. Do not add a fourth
% item — the page is full.
% GRADUAL RELEASE: the warm-up activates prior knowledge before the guided notes. All three are
% Algebra 1 tools the notes pick up again: item 1 (a percent from a part and a whole) is the
% summary a CATEGORICAL variable gets — problem 6 (50% captains); item 2 (scaling 18/50 up to
% 900) is sample -> population, which Lesson 1.2 formalizes and homework item 5 spirals; item 3
% (a mean, then a sentence naming the group and the units) is the "a number needs a context"
% idea row 1 turns on — problem 1 asks for exactly this sentence again.
\documentclass[12pt]{article}
\usepackage{saar-article}
\usepackage{saar-boxes}
\newcommand{\TallMath}[1]{$\displaystyle #1\rule[-1.4em]{0pt}{3.2em}$}
% Handwriting room inside every \begin{work} block. Moves the blank and the
% key together, so the two can never drift apart.
\setlength{\workrowsep}{50pt}

\begin{document}

\pageheader{Unit 1, Lesson 1.0}{Warm-Up}
% NAMESTRIP: no \namedateperiod here. The student writes their name once, on the
% cover this component is stapled behind.

\vspace{0.15in}

\begin{notesbox}{Spiral Review --- Percents, Proportions, and the Mean}
\normalsize
Everything below comes back this unit. Show your work.

\begin{enumerate}[leftmargin=1.8em, itemsep=22pt, topsep=10pt]

  \item Riverbend High School has $240$ seniors. Of those seniors, $84$ ride the
        bus to school. What \textbf{percent} of the seniors ride the bus?

        \begin{work}
          \dfrac{84}{240} &= 0.35 = 35\%
        \end{work}

  \item A survey of $50$ Riverbend students found that $18$ of them walk to school.
        Riverbend has $900$ students in all. Based on this survey, \textbf{predict}
        how many of the $900$ students walk to school.

        \begin{work}
          \dfrac{18}{50} &= 0.36 \\
          0.36 \times 900 &= 324 \text{ students}
        \end{work}

  \item Five students recorded how many minutes they spent on homework last night:
        $12$, $15$, $9$, $20$, $14$.

        \textbf{(a)} Find the \textbf{mean} number of minutes.

        \begin{work}
          \bar{x} &= \dfrac{12+15+9+20+14}{5} = \dfrac{70}{5} = 14 \text{ minutes}
        \end{work}

        \textbf{(b)} Write one sentence that says what your answer in part (a)
        describes. Name the \emph{group} and the \emph{units}.
        \par\writeline

\end{enumerate}
\end{notesbox}

\end{document}
